% !TEX root = ../main.tex
%\section{Method}
%\label{sec:method}


% % ----- Brian's standing comments on this section (collected at section top) -----
% \bc{I suggest that we simply have two sections: Preliminaries and <Method Name>. At the very beginning of the methods section, we give a high-level overview and summarize what we will derive and explain in each of the sections.}

% \bc{As shown in Fig.~X, higher frequencies of latents emerge naturally throughout the denoising process. To fully leverage this spectral autoregression property of diffusion models for massive computational savings, we propose \methodname, a novel dito seeffusion pipeline in which progressive resolution growth occurs within a single denoising trajectory via spectral transformations and noise padding. During denoising, partially denoised low-resolution latents are transformed to the spectral domain at closed-form transition times, augmented with high-frequency noise, and mapped back to the spatial domain via an inverse transform. This progressively increases resolution while avoiding computation on noise-dominated high-frequency tokens.}

% \bc{(on the power-spectrum paragraph, now in Preliminaries.) I would actually move this section to preliminaries. At the very beginning of Sec.~\ref{sec:method:noise-injection}, we can then clearly motivate our method with this subsection. See my rewrite of Sec.~\ref{sec:method:noise-injection}. The ``noise-floor'' plot where specific frequencies emerge should also happen before this section.}

% \bc{(on the training-free subsection.) I changed injection to expansion because this signals resolution growth.}
% % --------------------------------------------------------------------------------

\section{Preliminaries}
\label{sec:method:prelim}

\subsection{Diffusion Models and Flow Matching}
\label{sec:method:prelim:basics}

Diffusion models~\citep{ho2020denoising,song2021score} define a generative process that gradually transforms samples from a simple prior (e.g. Gaussian) to data samples via a learned reverse-time process. Flow Matching \cite{liu2022rectified,lipman2023flow} reformulates the diffusion process as a continuous-time optimal transport problem and learns a velocity field that deterministically transforms noisy samples to data samples along straight paths.

Specifically, given a clean data point $\mathbf{x}_0$ sampled from a data distribution $p_{\text{data}}$, and a noise sample $\mathbf{x}_1 = \boldsymbol{\epsilon} \sim \mathcal{N}(0, I)$, the noise-to-data path is a straight line defined by
\begin{equation}
  \mathbf{x}_t \;=\; (1-t)\, \mathbf{x}_0 \,+\, t\,\boldsymbol{\epsilon},
  \qquad t \in [0, 1].
  \label{eq:flow-forward}
\end{equation}
A neural network
$\mathbf{v}_\theta(\mathbf{x}_t, t)$ is trained to predict the target velocity $\dot{\mathbf{x}}_t = \boldsymbol{\epsilon} - \mathbf{x}_0$ with the following objective:
\begin{equation}
  \mathcal{L}_\text{flow}(\theta) \;=\; \mathbb{E}_{t, \mathbf{x}_0, \boldsymbol{\epsilon}}\!\left[\,\|\mathbf{v}_\theta(\mathbf{x}_t, t) - (\boldsymbol{\epsilon} - \mathbf{x}_0)\|^2\,\right].
  \label{eq:flow-loss}
\end{equation}

A network with sufficient capacity that is trained on enough data converges to the \textit{Bayes-optimal velocity predictor} $\mathbf{v}_\theta(\mathbf{x}_t, t) = \mathbf{v}^*(\mathbf{x}_t, t)$~\citep{liu2022rectified, lipman2023flow}.

\iffalse

Under the flow-matching objective in Eq.~\eqref{eq:flow-loss}, the Bayes-optimal velocity predictor is the conditional expectation
\begin{equation}
  \mathbf{v}^*(\mathbf{x}_t, t) \;=\; \mathbb{E}\!\left[\,\boldsymbol{\epsilon} - \mathbf{x}_0 \,\big|\, \mathbf{x}_t\,\right],
  \label{eq:optimal-v}
\end{equation}
which is the unique minimizer of Eq.~\eqref{eq:flow-loss}~\citep{liu2022rectified,lipman2023flow}. A network with sufficient capacity trained on enough data converges to $\mathbf{v}^*$ in expectation.
\fi


During inference, data samples are generated through sampling from the prior $\mathbf{x}_1 \sim \mathcal{N}(0, I)$ and solving the probability-flow Ordinary Differential Equation (ODE) $\dot{\mathbf{x}}_t = \mathbf{v}_\theta(\mathbf{x}_t, t)$. Due to its faster training and stable convergence, recent image and video generative models \cite{flux-2,zimage,sd3,wan} have increasingly adopted the flow matching paradigm.


Most state-of-the-art diffusion and flow matching models adopt the Diffusion Transformer (DiT) architecture \citep{dit}, which processes tokens through a sequence of self-attention and Multilayer Perceptron (MLP)-based operations. This motivates our approach to reduce the number of tokens processed along the denoising trajectory using the spectral autoregression property of diffusion models \citep{dieleman2024spectral}.

%costing one forward pass of $v_\theta$ per step. 

%Our framework applies to both diffusion \eqref{eq:diff-forward} and flow-matching \eqref{eq:flow-forward} backbones.


% Modern image and video diffusion models often operate in a compressed latent space~\citep{rombach2022ldm} rather than on raw pixels.

\iffalse
Diffusion models~\citep{ho2020denoising,song2021score} define a forward process that progressively corrupts a clean sample $x_0 \sim p_\text{data}$ by adding i.i.d.\ Gaussian noise,
\begin{equation}
  x_t \;=\; \alpha_t\, x_0 \,+\, \sigma_t\,\epsilon,
  \qquad t \in [0, 1], \qquad \epsilon \sim \mathcal{N}(0, I),
  \label{eq:diff-forward}
\end{equation}
with schedule $(\alpha_t, \sigma_t)$ moving from $(1, 0)$ at $t = 0$ (clean data) to $(0, 1)$ at $t = 1$ (pure noise). A network $\epsilon_\theta(x_t, t)$ is trained to predict the noise by minimizing
\begin{equation}
  \mathcal{L}_\text{diff}(\theta) \;=\; \mathbb{E}_{t, x_0, \epsilon}\!\left[\,\|\epsilon_\theta(x_t, t) - \epsilon\|^2\,\right].
  \label{eq:diff-loss}
\end{equation}
Generation starts from a random sample $x_1 = \epsilon \sim \mathcal{N}(0, I)$ and numerically integrates the corresponding reverse-time process (the probability-flow ODE) backward from $t = 1$ to $t = 0$, with the score $\nabla_{x_t}\!\log p_t(x_t) = -\epsilon_\theta(x_t, t)/\sigma_t$ obtained via Tweedie's identity~\citep{efron2011tweedie}. 
\fi

\iffalse
\paragraph{Diffusion Transformers.}
\bc{can probably remove this section, since we will explain this in intro already}
Modern image and video generation models are predominantly built on Diffusion Transformer (DiT) backbones~\citep{dit}, where image and video pixels or latents are encoded into a sequence of $N$ tokens, then processed by a stack of self-attention and feed-forward MLP blocks. In particular, the computational cost of self-attention is quadratic in $N$, making efficient, high-resolution generation especially challenging. This motivates our approach: by restricting early denoising steps where high-frequency components remain noise-dominated to a reduced spatial resolution, we avoid unnecessary computation on uninformative high-frequency tokens, thereby significantly improving generation efficiency.
\fi

% Let $\mathbf{x}_0 \in \mathbb{R}^{C \times T \times H \times W}$ denote a clean data sample drawn from $p_\text{data}$ ($C$ channels and $T$ frames, where $T = 1$ for images and $T > 1$ for videos).


% !TEX root = ../main.tex
\begin{figure}[t]
  \centering
  \includegraphics[width=\linewidth]{figures/power_spectrum.pdf}
  \caption{\textbf{Diffusion process in the spectral domain.} Latent power spectra in both image and video models decay rapidly with frequency (Fig.~\textbf{(a)}), consistent with natural images. Diffusion exhibits a frequency-domain autoregressive structure (Fig.~\textbf{(b)}) due to the aforementioned property: low frequencies emerge early in the denoising process, while high frequencies remain noise-dominated.}
  \label{fig:wavelet-spectral}
\end{figure}


\subsection{Spectral Autoregression in Diffusion Models}
\label{sec:method:prelim:spectral}

Natural images exhibit a characteristic power-law decay \citep{ruderman1997origins}, which is the foundation of the spectral autoregression property of diffusion models \citep{dieleman2024spectral}. Specifically, let $\Phi = \{\phi_\omega\}_{\omega \in \Omega}$ be a spectral basis indexed by frequency $\omega \in \Omega$ (e.g., Fourier~\citep{fourier1822theorie}, discrete cosine~\citep{ahmed1974dct}, or Haar wavelet~\citep{haar1909theorie}), and denote the associated spectral transformation by $T_\Phi$. We define the per-frequency signal power, or \emph{power spectrum}, of a data distribution $p_{\text{data}}$ as
\begin{equation}
  P_\omega \;:=\; \mathbb{E}_{\mathbf{x}_0}\!\left[\,\bigl|x_0^{(\omega)}\bigr|^2\,\right],
  \quad \text{where  } x_0^{(\omega)} := \langle \mathbf{x}_0, \phi_\omega \rangle \text{ under $T_\Phi$}.
  \label{eq:Pb-def}
\end{equation}

When $p_\text{data}$ represents natural images, $P_\omega$ is known to exhibit a power-law decay, that is
\begin{equation}
  P_\omega \;\propto\; |\omega|^{-\beta},
  \label{eq:power-law}
\end{equation}
with the exponent term typically in the range $\beta \in [2, 3]$~\citep{dieleman2024spectral,fourierspace}. We validated similar trends in both image and video latent representations of modern latent-space diffusion models (see Fig. \ref{fig:wavelet-spectral}).

% Although modern diffusion models operate in VAE latent space rather than pixel space \citep{rombach2022ldm}, their frequency statistics remain largely preserved. We compute radially averaged power spectra of latent representations of image and video generation models and observe a consistent power-law decay (Fig. \ref{fig:wavelet-spectral}). 

%We estimate exponents of $\beta \approx 1.92$ for FLUX.1-dev image latents \citep{black2024flux} and $\beta \approx 2.42$ for WAN 2.1 video latents \citep{wan} (Fig.~\ref{fig:wavelet-spectral}).

The power-law decay of $P_\omega$ in Eq.~\eqref{eq:power-law} implies that high frequencies carry substantially weaker signal than low frequencies. Since the Gaussian noise $\boldsymbol{\epsilon}$ added during the forward process has a flat, i.e. frequency-independent, power spectrum, higher frequencies are dominated by noise early in the denoising trajectory. In this work, we leverage this spectral autoregression property to progressively grow resolution within a single denoising trajectory via spectral domain transformations for efficient image and video generation.

% As observed by \citet{dieleman2024spectral} and \citet{fourierspace}, diffusion models exhibit the spectral autoregression property, recovering low-frequency signals early in the denoising process and generating higher-frequency details only in the later denoising steps.


\section{Spectral Progressive Diffusion}
In this section, we present a spectral-transformation-based progressive generation framework that enables arbitrary resolution transitions within a single denoising trajectory for pretrained image and video generation models (Sec.~\ref{sec:method:noise-injection}). We then derive an optimal resolution transition schedule based on power spectrum analysis and resolution-dependent frequency limits (Sec.~\ref{sec:method:schedule:signal}). Finally, we introduce a principled fine-tuning strategy to further improve both generation quality and efficiency (Sec.~\ref{sec:method:lora}).

\label{sec:method:schedule}

% ----- Commented-out intro (previous 3.2 opening paragraph), kept for reference: -----
% This section develops the theoretical foundation of our resolution schedule. We project the latent onto an orthonormal basis of the frequency domain, decomposing flow matching into a collection of independent per-frequency subproblems, and show that the expected velocity prediction error on each frequency is determined entirely by the signal's power spectrum. This yields a principled per-frequency activation criterion based on when the velocity is reliably predictable from the noisy input. In this section we include our derivations for the flow-matching parameterization of Eq.~\eqref{eq:flow-forward}; the analogous derivation for diffusion models is provided in the supplementary material.
% ------------------------------------------------------------------------------------


\subsection{Training-Free Inference with Spectral Noise Expansion}
\label{sec:method:noise-injection}

% In practice, Eq.~\eqref{eq:stage-transition} is discretized into $S$ resolution stages $(0 < s_1 < s_2 < \cdots < s_S = 1)$, with stage $i$ held for $t \in (t_{s_i}^{\text{transition}},\, t_{s_{i-1}}^{\text{transition}}]$ and matched to the denoising timesteps of the pretrained models.

Motivated by the spectral autoregression property of diffusion (Sec.~\ref{sec:method:prelim:spectral}), we progressively increase image resolution by injecting higher-frequency components along the denoising trajectory using a spectral transformation. Our framework applies natively to existing pretrained image and video generation models, without additional custom modules or architectural modifications.

Given an orthonormal spectral basis $\Phi$, let $T_\Phi$ and $T_\Phi^{-1}$ denote the forward and inverse spectral transformations. We define $S$ progressive resolution scales $s_{1:S}$, where $0 < s_1 < s_2 < \cdots < s_S = 1$, and their corresponding resolution transition times $t_{1:S-1}$, with $1 > t_1 > \cdots > t_{S-1} > 0$, that are matched to the noise schedules of pretrained models. We denote the final full-resolution state as $\mathbf{x}_0^{s_S} \in \mathbb{R}^{C \times T \times H \times W}$, and assume stage $i$ runs at spatial resolution $(s_iH, s_iW)$ for $t \in (t_i,\, t_{i-1}]$. 
% In practice, we select $s_i$ such that resolutions $(s_i H, s_i W)$ align with the multi-resolution training distribution of pretrained generative models. For example, we set $s_i$ to be powers of $1/2$ for FLUX.1-dev, as it is primarily trained on image latents with power-of-two resolutions.

At each transition time $t_i$, we expand the current low-resolution state $\mathbf{x}^{s_i}_{t_i} \in \mathbb{R}^{C \times T \times (s_iH) \times (s_iW)}$ to the next resolution state $\mathbf{x}^{s_{i+1}}_{t_i} \in \mathbb{R}^{C \times T \times (s_{i+1}H) \times (s_{i+1} W)}$ using \emph{spectral noise expansion} and timestep alignment (see Fig.~\ref{fig:method}).

\paragraph{Spectral noise expansion.}
Given a low-resolution $\mathbf{x}^{s_i}_{t_i}$ at transition time $t_i$, we expand its resolution to the next scale $s_{i + 1}$ through the following steps:
\begin{enumerate}
  \item[(i)] Compute the spectrum $\boldsymbol{\xi}^{s_i}_{t_i} = T_\Phi(\mathbf{x}^{s_i}_{t_i})$ supported on the frequency set $\Omega_{s_i}$ representable at scale $s_i$.
  \item[(ii)] Embed $\boldsymbol{\xi}^{s_i}_{t_i}$ in the lower-frequency part of the spectrum for $\Omega_{s_{i + 1}}$ and fill the slots $\Omega_{s_{i+1}} \setminus \Omega_{s_i}$ with $t_i \epsilon$, $\epsilon \sim \mathcal{N}(0, 1)$ to get $\boldsymbol{\xi}^{s_{i + 1}}_{t_i}$. 
  \item[(iii)] Convert the spectrum back to spatial domain $\mathbf{x}^{s_{i+1}}_{t_i}$ at resolution $(s_{i+1}H, s_{i+1}W)$ via $\mathbf{x}^{s_{i+1}}_{t_i} = T_\Phi^{-1}(\boldsymbol{\xi}^{s_{i + 1}}_{t_i})$.
\end{enumerate}
Steps~(i)--(iii) expand the resolution scale from $s_i$ to $s_{i + 1}$ for $\mathbf{x}_{t_i}^{s_{i + 1}}$ by injecting high frequencies at the correct noise level while preserving the partially-denoised low-frequency content of $\mathbf{x}^{s_i}_{t_i}$.

\paragraph{Timestep alignment.}
As $\Phi$ is orthonormal, the high-frequency noise padding in $\Omega_{s_{i + 1}} \setminus \Omega_{s_i}$ raises the overall noise level of the enlarged $\mathbf{x}_{t_i}^{{s_{i+1}}}$, which makes it no longer correspond to the original timestep $t_i$. We therefore rescale the inverse-transformed output and re-index the timestep by:
\begin{equation}
  \tilde{\mathbf{x}}^{s_{i+1}}_{\tilde{t}_i} := \frac{s_{i+1}/s_i}{1 + ((s_{i+1}/s_i) - 1)t_i} \cdot \mathbf{x}^{s_{i+1}}_{t_i},
  % \tilde{\mathbf{x}}^{s_{i+1}}_{\tilde{t}_i} = \frac{s_{i+1}/s_i}{1 + ((s_{i+1}/s_i) - 1)t_i} \cdot \mathbf{x}^{s_{i+1}}_{t_i},
  \label{eq:scale-align}
\end{equation}
where
\begin{equation}
  \tilde{t}_i := \frac{(s_{i+1}/s_i)\,t_i}{1 + ((s_{i+1}/s_i) - 1)t_i}.
  \label{eq:time-align}
\end{equation}
Eq.~\ref{eq:scale-align} and Eq.~\ref{eq:time-align} give us the aligned state $\tilde{\mathbf{x}}^{s_{i+1}}_{\tilde{t}_i}$ at the aligned timestep $\tilde{t}_i$. The integration of the probability-flow ODE then resumes at resolution $(s_{i+1}H, s_{i+1}W)$ until the next transition time $t_{i+1}$. 
The full derivation of Eqs.~\eqref{eq:scale-align}--\eqref{eq:time-align} is provided in Appendix~\ref{app:time-align}. In practice, $S$ and $s_{1:S}$ are chosen to align with the multi-resolution training distribution of pretrained generative models. By default, we choose $T_\Phi$ as the Discrete Cosine Transform (DCT)~\citep{ahmed1974dct} with discrete cosine basis $\Phi$. Alternative spectral transforms (such as Discrete Wavelet Transform (DWT) and Fourier Transform (FFT)) are compared and evaluated in Sec.~\ref{sec:exp:ablations}.


\renewcommand{\arraystretch}{1.0}
\setlength{\aboverulesep}{0.2ex}
\begin{table*}[!t]
  \centering
  \scriptsize
  \begin{tabular*}{\textwidth}{@{\extracolsep{\fill}} l ccccccc @{}}
    \toprule
    \multirow{2}{*}{Method} & \multirow{2}{*}{\makecell{Speedup (s) $\uparrow$}} & \multirow{2}{*}{\makecell{TFLOPs $\downarrow$}} & Overall & \multicolumn{2}{c}{Image quality} & \multicolumn{2}{c}{Text alignment} \\
    \cmidrule(lr){4-4} \cmidrule(lr){5-6} \cmidrule(lr){7-8}
     & & & ImageReward $\uparrow$ & CLIP-IQA $\uparrow$ & NIQE $\downarrow$ & T2I-Comp. $\uparrow$ & GenEval $\uparrow$ \\
    \midrule
    FLUX (50 steps)             & 1.00$\times$              & 2991.01              & \textbf{1.095}    & 0.707             & 6.75             & \textbf{0.634}    & \textbf{0.698} \\
    RALU~\citep{ralu}     & \underline{1.58}$\times$              & \textbf{1749.94}  & 1.028             & \underline{0.712} & \textbf{6.07}    & 0.613                 & 0.648     \\
    \textbf{Ours ($S=2$)} & \textbf{1.66}$\times$  & \underline{1755.22}              & \underline{1.049} & \textbf{0.719}    & \underline{6.43}             & \underline{0.617} & \underline{0.654}     \\
    % \textbf{Ours ($S=3$)} & \textbf{1.73}$\times$     & \textbf{1.79}$\times$     & 1.015             & 0.711             & \underline{6.33} & 0.593             & 0.640     \\
    \midrule[0.3pt]
    FLUX (10 steps) & 4.84$\times$ & 610.02 & 0.981 & 0.679 & 6.93 & 0.618 & 0.647 \\
    Bottleneck~\citep{bottleneck} & 4.67$\times$              & 571.23              & 0.889          & 0.661             & 9.16          & 0.620             & \textbf{0.687}    \\
    RALU~\citep{ralu}             & 4.98$\times$              & 540.47              & 1.022          & \underline{0.700} & \textbf{6.43} & \textbf{0.626}    & 0.652             \\
    \textbf{Ours ($S=2$)}         & \underline{5.77}$\times$  & \underline{500.34}  & \textbf{1.059} & 0.696             & 6.69          & \underline{0.624} & \underline{0.655} \\
    \textbf{Ours ($S=3$)}         & \textbf{6.09}$\times$     & \textbf{469.15}     & \underline{1.042} & \textbf{0.701} & \underline{6.53} & 0.623          & 0.637             \\
    \midrule[0.3pt]
    FLUX (7 steps) & 6.62$\times$ & 431.45 & 0.920 & 0.660 & 8.25 & 0.594 & 0.583 \\
    Bottleneck~\citep{bottleneck} & 6.64$\times$              & 431.52              & 0.792          & 0.631             & 8.71          & 0.605             & \underline{0.672} \\
    RALU~\citep{ralu}             & 6.69$\times$              & \underline{426.01}  & 0.999          & 0.681             & 6.87          & \textbf{0.633}    & \textbf{0.682}    \\
    \textbf{Ours ($S=2$)}         & \underline{6.78}$\times$  & 427.03              & \textbf{1.039} & \underline{0.689} & \underline{6.78} & 0.620          & 0.667             \\
    \textbf{Ours ($S=3$)}         & \textbf{7.09}$\times$     & \textbf{406.24}     & \underline{1.015} & \textbf{0.694} & \textbf{5.99} & \underline{0.627} & 0.637             \\
    \bottomrule
  \end{tabular*}
  \vspace{-2pt} 
  \caption{\textbf{Training-free quantitative comparisons on FLUX.1-dev~\citep{black2024flux}.} Image resolution is $1024^2$ and baseline rows are copied from Table 1 in RALU~\citep{ralu}. We group results based on $1\times, 5\times, 7\times$ speedup settings where progressive resolution approaches are applied in combination with reducing denoising steps. Our method improves the speed--quality tradeoff across all speedup settings compared to baseline approaches. Further evaluation details and comparisons are shown in Appendix Sec.~\ref{subsec:image_latent}.}
  % replicating RALU's~\citep{ralu} Table~1 protocol; we keep only spatial-acceleration baselines.
  \label{tab:flux-main}
\end{table*}



\subsection{Optimal Resolution Schedule}

\label{sec:method:schedule:signal}

Section~\ref{sec:method:noise-injection} described our progressive resolution transition mechanism during the denoising process given a predetermined resolution transition schedule $t_{1:S-1}$. In practice, while increasing the number of resolution scales $S$ improves token efficiency, it introduces additional transition times as hyperparameters to tune. As a result, existing progressive resolution approaches typically limit $S = 2$ or leave $t_{1:S-1}$ as a brittle, model-specific design choice~\citep{bottleneck, ralu}. 
In contrast, our method supports an arbitrary number of resolutions $S$ and arbitrary resolution scales $s_{1:S}$. In the following section, we derive a $\delta$-\textit{optimal resolution schedule} for pretrained flow matching models that determines the optimal transition times $t^*_{1:S-1}$ from a single error threshold parameter $\delta$ independent of $S$.

%We include similar derivations for standard diffusion models in the supplemental materials.

\paragraph{Spectral-domain flow matching.}
The spatial-domain flow matching forward process defined in Eq.~\eqref{eq:flow-forward} can be written in the spectral domain under transformation $T_\Phi$ with orthonormal basis $\Phi = \{\phi_\omega\}_{\omega \in \Omega}$:
\begin{equation}
  x_t^{(\omega)} \;=\; (1-t)\,x_0^{(\omega)} \,+\, t\,\epsilon^{(\omega)},
  \qquad \omega \in \Omega,
  \label{eq:forward-mode}
\end{equation}
where $\epsilon^{(\omega)} := \langle \boldsymbol{\epsilon}, \phi_\omega \rangle$. By orthonormality of $\Phi$, $\epsilon^{(\omega)} \sim \mathcal{N}(0, 1)$ is i.i.d.\ Gaussian noise for each $\omega \in \Omega$. We denote the per-frequency component of the Bayes-optimal velocity predictor described in Sec. \ref{sec:method:prelim:basics} by $v^{*(\omega)}(\mathbf{x}_t, t)$. We then determine when a spectral component $\omega$ is noise-dominated in Proposition 1 and link this to spatial resolution to derive the $\delta$-optimal resolution schedule in Proposition 2. 

% We can then decompose the flow matching training objective defined in Eq.~\eqref{eq:flow-loss} into the sum of individual objectives for each $\omega \in \Omega$ using Parseval's identity~\citep{mallat1989wavelet}:
% \begin{equation}
%   \mathcal{L}_\text{flow}(\theta) \;=\; \sum_{\omega \in \Omega}\, \mathbb{E}_{t,\mathbf{x}_0,\mathbf{\epsilon}}\!\left[\,\bigl|v_\theta^{(\omega)}(\mathbf{x}_t, t) - (\epsilon^{(\omega)} - x_0^{(\omega)})\bigr|^2\,\right].
% \end{equation}

\textbf{Proposition 1 (Per-frequency $\delta$-optimal activation time):} Under the simplified modelling assumption $x_0^{(\omega)} \sim \mathcal{N}(0, P_\omega)$ with $P_\omega > 0$, for all $t \ge t_\omega$ and $\delta\in(0, 1)$, we have:
\begin{equation}
  \mathbb{E}\!\left[\bigl|v^{*(\omega)}(\mathbf{x}_t, t) - \epsilon^{(\omega)}\bigr|^2\right] \;\le\; \delta,
  \label{eq:noise-dominated}
\end{equation}
where $t_\omega$ is the $\delta$-optimal activation time for $\omega$:
\begin{equation}
  t_\omega \;:=\; \frac{1}{1 + \sqrt{\dfrac{\delta}{P_\omega\,(1 + P_\omega - \delta)}}}.
  \label{eq:activation}
\end{equation}
Intuitively, for $t \geq t_\omega$, spectral component $\omega$ is noise-dominated as its optimal velocity predictions are approximately noise $\epsilon^{(\omega)}$. After $t_\omega$, the signal power overcomes the error tolerance $\delta$ and structural information is recovered. The proof of Proposition 1 is given in Appendix Sec.~\ref{app:proof-prop1}. We further include empirical evidence supporting Proposition 1 in Appendix Sec.~\ref{sec:motivation}.


\textbf{Proposition 2 (Per-resolution $\delta$-optimal transition time):} Under the setting of Proposition~1, and assuming $P_\omega$ is monotonically decreasing in $|\omega|$ (consistent with the power-law decay Eq.~\eqref{eq:power-law}), for any resolution scale $s_i, s_{i + 1}$, the optimal transition time from scale $s_i$ up to $s_{i+1}$ is
\begin{equation}
  t_i^* \;:=\; \min_{\omega \in \Omega_{s_i}}\, t_\omega \;=\; t_{\omega=s_i \cdot \omega_{\max}(H,W)},
  \label{eq:stage-transition}
\end{equation}
where $\Omega_{s_i}$ is the set of frequencies representable on the $(s_iH, s_iW)$ grid. The maximum representable frequency of the full-resolution spatial grid $\Omega_{s_S} = \omega_{\max}(H, W) = \min(H, W)/2$ is given by the Nyquist--Shannon sampling theorem~\citep{nyquist1928certain,shannon1949communication}. 

Proposition 2 connects frequency activation times with spatial resolution transitions through the Nyquist limit. It implies that high-resolution computation in early denoising steps is redundant as most representable frequencies $\omega$ are still noise-dominated, motivating our approach of Spectral Progressive Diffusion. We apply spectral noise expansion at $t_i^*$ to progressively increase resolution precisely when finer details emerge from noise, thereby maximizing token efficiency. The hyperparameter $\delta$ is independent of $S$, making our method particularly robust to tuning compared to prior methods. The proof of Proposition 2 is deferred to Appendix~\ref{app:proof-prop2} and we analyze the effect of the number of resolution stages $S$ and the tolerance $\delta$ in Sec.~\ref{sec:exp:ablations}.

% In practice, Eq.~\eqref{eq:stage-transition} is discretized into $S$ resolution stages $(0 < s_1 < s_2 < \cdots < s_S = 1)$, with stage $i$ held for $t \in (t_{s_i}^{\text{transition}},\, t_{s_{i-1}}^{\text{transition}}]$ and matched to the denoising timesteps of the pretrained models.

% \paragraph{Per-frequency SNR.}
% Under the basis $\Phi$, the rectified-flow forward process in Eq.~\eqref{eq:flow-forward} factors into independent per-frequency equations in which the noise component $\epsilon^{(\omega)} \sim \mathcal{N}(0, 1)$ is i.i.d.\ Gaussian (full derivation in Sec.~\ref{sec:method:schedule:signal}). Following~\citet{kingma2021vdm}, we define the \emph{per-frequency signal-to-noise ratio} as the ratio of the clean- and noise-component second moments in this process:
% \begin{equation}
%   \mathrm{SNR}_\omega(t)
%   \;:=\;
%   \frac{\mathbb{E}[\,|(1-t)\,x_0^{(\omega)}|^2\,]}{\mathbb{E}[\,|t\,\epsilon^{(\omega)}|^2\,]}
%   \;=\;
%   \underbrace{\frac{(1-t)^2}{t^2}}_{\mathrm{SNR}(t)}\;\cdot\;P_\omega,
%   \label{eq:snr}
% \end{equation}
% which factorises into the frequency-independent timestep-SNR term $\mathrm{SNR}(t) = (1-t)^2/t^2$ and the frequency-specific signal power $P_\omega$.

% \paragraph{Frequency limit from resolution.}
% By the Nyquist--Shannon sampling theorem~\citep{nyquist1928certain,shannon1949communication}, when generating on a spatial grid of size $H \times W$ the representable 2D frequencies form the rectangle $|\omega_x| \le W/2$, $|\omega_y| \le H/2$. For the radial analysis used to derive our resolution schedule, we define the maximum representable radial frequency as
% \begin{equation}
%   \omega_{\max}(H, W) \;=\; \frac{\min(H, W)}{2}.
%   \label{eq:nyquist-cap}
% \end{equation}
% A stage running at reduced resolution $(sH, sW) \le (H, W)$ for $s \leq 1$ can represent only the frequency subset
% \begin{equation}
%   \Omega_s \;:=\; \{\,\omega \in \Omega : |\omega| \le s \cdot \omega_{\max}(H, W)\,\} \;\subseteq\; \Omega.
%   \label{eq:nyquist-set}
% \end{equation}



% To characterize how much useful signal each frequency carries at time $t$, we measure how far $v^*$ deviates from the null prediction $\hat v = \epsilon$. Under the simplified modelling assumption that $x_0^{(\omega)} \sim \mathcal{N}(0, P_\omega)$, we get:
% \begin{equation}
%   \mathbb{E}\!\left[\bigl|v^{*(\omega)} - \epsilon^{(\omega)}\bigr|^2\right]
%   \;=\; \frac{\mathrm{SNR}_\omega(t)\,(1 + P_\omega)}{1 + \mathrm{SNR}_\omega(t)},
%   \label{eq:velocity-error-optimal}
% \end{equation}
% We defer the derivations of above to Appendix~\ref{app:bayes-velocity}. 

% \paragraph{Noise-dominated frequencies.}
% Given a tolerance $\delta > 0$, we say frequency $\omega$ is \emph{noise-dominated at time $t$} if the Bayes-optimal velocity deviates from the null prediction by at most $\delta$ in expectation:
% \begin{equation}
%   \mathbb{E}\!\left[\bigl|v^{*(\omega)} - \epsilon^{(\omega)}\bigr|^2\right] \;\le\; \delta.
%   \label{eq:noise-dominated}
% \end{equation}
% Substituting the Bayes-optimal error of Eq.~\eqref{eq:velocity-error-optimal} and the SNR factorisation from Eq.~\eqref{eq:snr} into this definition, and solving for $t$, turns the noise-dominated condition into an inequality on the timestep:
% \begin{equation}
%   t \;\ge\; \frac{1}{1 + \sqrt{\dfrac{\delta}{P_\omega\,(1 + P_\omega - \delta)}}}
%   \;=:\; t_\omega,
%   \label{eq:activation}
% \end{equation}
% defining the frequency-specific \emph{activation time} $t_\omega$ as the right-hand side. We assume $\delta < 1$, and the right hand side of Eq.~\eqref{eq:activation} is well defined with $0 < t_\omega < 1$. Frequency $\omega$ is thus noise-dominated for $t \ge t_\omega$ and transitions to signal-carrying for $t < t_\omega$ under tolerance $\delta$.

% \paragraph{$\delta$-optimum resolution schedule.}
% A \emph{resolution schedule} assigns each timestep $t$ a resolution scale $s_t \in (0, 1]$ such that the generator runs on the grid $(s_t H,\, s_t W)$ and represents the frequencies in $\Omega_{s_t}$ of Eq.~\eqref{eq:nyquist-set}. We call a schedule \emph{$\delta$-optimum} if at every $t$ it uses the smallest scale that still represents every frequency already signal-carrying under tolerance $\delta$. Formally,
% \begin{equation}
%   s^*_t \;=\; \min\,\bigl\{\,s \in (0, 1] : \{\omega \in \Omega : t_\omega > t\} \subseteq \Omega_s\,\bigr\}
%   \;=\; \frac{\omega^*(t)}{\omega_{\max}(H, W)},
%   \label{eq:delta-opt-schedule}
% \end{equation}
% where $\omega^*(t) := \max\{|\omega| : t_\omega > t\}$ is the highest frequency that has crossed its activation time $t_\omega$ (i.e., frequencies with $t < t_\omega$ are already signal-carrying at time $t$).

% Equivalently, the transition time from resolution scale $s$ up to a finer resolution scale $s' > s$ is exactly when every $\omega \in \Omega_s$ has become signal-carrying, i.e., when $t$ falls to the smallest activation time within $\Omega_s$:
% \begin{equation}
%   t_s^{\text{transition}} \;:=\; \min_{\omega \in \Omega_s}\, t_\omega \;=\; t_{\,s \cdot \omega_{\max}(H,W)},
%   \label{eq:stage-transition}
% \end{equation}
% where the second equality follows from the fact that $t_\omega$ is decreasing in $|\omega|$ as $P_\omega$ decays with $|\omega|$ (Eq.~\eqref{eq:power-law}), and Eq.~\eqref{eq:activation} is monotonically increasing in $P_\omega$.

\renewcommand{\arraystretch}{1.0}
\setlength{\doublerulesep}{0.5pt}
\setlength{\aboverulesep}{0.2ex}
\begin{table*}[!t]
  \centering
  \scriptsize
  \begin{tabular*}{\textwidth}{@{\extracolsep{\fill}} l ccccccc @{}}
    \toprule
    Method & \makecell{Speedup (s) $\uparrow$} & \makecell{TFLOPs $\downarrow$} & ImageReward $\uparrow$ & CLIP-IQA $\uparrow$ & NIQE $\downarrow$ & T2I-Comp. $\uparrow$ & GenEval $\uparrow$ \\
    \midrule
    Z-Image (50 steps)           & 1.00$\times$             & 4941.23             & 0.965    & 0.700    & 5.41    & 0.731    & 0.745 \\
    % Z-Image (32 steps)           & {1.56}$\times$             & {3166.62}             & 0.957    & 0.697    & {5.44}    & 0.686    & 0.725 \\
    \midrule
    \textbf{Ours} (TF, $S=2$)  & {1.65}$\times$ & {3132.03} & {0.904}             & {0.688}             & {5.87}             & {0.658} & {0.730} \\
    % \textbf{Ours ($S=3$)}  & \textbf{1.74}$\times$    & \textbf{1.72}$\times$    & 0.875             & 0.690             & {5.59} & 0.650             & 0.682 \\
    \textbf{Ours} (LoRA, $S = 2$)   & {1.65}$\times$             & {3132.03}             & \textbf{0.982} & \textbf{0.699} & \textbf{5.72}             & \textbf{0.725}                 & \textbf{0.731} \\
    \midrule
    \textbf{Ours} (TF, $S=3$)  & {1.74}$\times$ & {2871.09} & {0.875}             & {0.690}             & \textbf{5.59}             & {0.650} & {0.682} \\
    % \textbf{Ours ($S=3$)}  & \textbf{1.74}$\times$    & \textbf{1.72}$\times$    & 0.875             & 0.690             & {5.59} & 0.650             & 0.682 \\
    \textbf{Ours} (LoRA, $S = 3$)   & {1.74}$\times$             & {2871.09}             & \textbf{0.954} & \textbf{0.697} & {5.75}             & \textbf{0.717}                 & \textbf{0.728} \\
    % \textbf{Ours} (LoRA, $\delta=0.05$)   & \textbf{2.01}$\times$             & \textbf{2436.19}             & {0.919} & 0.684 & 6.12             & {0.661}                 & 0.718 \\
    % \midrule
    % \textbf{Ours} (TF, $S=3$)  & \textbf{1.74}$\times$ & \textbf{2871.09} & 0.875            & 0.690            & 5.59             & 0.650 & 0.682 \\
    % \textbf{Ours} (LoRA, $S=3$)  & \textbf{1.74}$\times$ & \textbf{2871.09} & \textbf{0.954}             &  \textbf{0.697}            & 5.75             & - & - \\

    % \textbf{Ours}
    \midrule
    \midrule
    % \textbf{Ours} (TF, $S=3$)  & \textbf{1.74}$\times$ & \textbf{2871.09} & 0.904             & 0.688             & 5.87             & {0.658} & {0.730} \\
    % % \textbf{Ours ($S=3$)}  & \textbf{1.74}$\times$    & \textbf{1.72}$\times$    & 0.875             & 0.690             & {5.59} & 0.650             & 0.682 \\
    % \textbf{Ours} (LoRA, $S=3$)   & \textbf{1.74}$\times$             & \textbf{2871.09}             & \textbf{0.982} & {0.699} & \textbf{5.72}             & \hx{TODO}                 & \hx{TODO} \\
    PixelGen (25 steps)         & 1.00$\times$             & 65.36             & 0.921    & 0.734    & 5.95    & 0.574    & 0.794 \\
    % PixelGen (13 steps)         & \textbf{1.92}$\times$             & {34.16}             & 0.886    & 0.726    & {5.89}    & 0.571    & 0.781 \\
    \midrule
    \textbf{Ours} (TF, $S=2$) & {1.60}$\times$             & {33.72}             & {0.799}    & {0.718}    & {6.10}    &  {0.568}  & \textbf{0.782} \\
    \textbf{Ours} (LoRA, $S=2$) & {1.55}$\times$             & {33.72}             & \textbf{0.913}    & \textbf{0.728}    & \textbf{5.87}    & \textbf{0.580}    & {0.776} \\
    \bottomrule
  \end{tabular*}
  % \vspace{3pt}
  \caption{\textbf{Fine-tuning quantitative comparisons on latent- (Z-Image~\citep{zimage}) and pixel-space image generation (PixelGen~\citep{pixelgen}).} Image resolution is $1024^2$ for Z-Image and $512^2$ for PixelGen. Across both latent- and pixel-space image generation, our fine-tuning method bridges the gap between model pretraining and progressive resolution inference, further improving generation quality while preserving efficiency gains. Further evaluation details and comparisons are shown in Appendix Sec.~\ref{subsec:image_pixel}.}
  \label{tab:zimage-main}
\end{table*}

\renewcommand{\arraystretch}{1.0}
\setlength{\aboverulesep}{0.2ex}
\begin{table*}[!t]
  \centering
  \scriptsize
  \begin{tabular*}{\textwidth}{@{\extracolsep{\fill}} l cccccccc @{}}
    \toprule
    \makecell{Method} &
    \makecell{Speedup (s) $\uparrow$} &
    \makecell{TFLOPs $\downarrow$} &
    \makecell{Subject\\Consistency $\uparrow$} &
    \makecell{Background\\Consistency $\uparrow$} &
    \makecell{Motion\\Smoothness $\uparrow$} &
    \makecell{Dynamic\\Degree $\uparrow$} &
    \makecell{Aesthetic\\Quality $\uparrow$} &
    \makecell{Image\\Quality $\uparrow$} \\
    \midrule
    WAN~2.1 (50)           & 1.00$\times$ & 119292         & 0.9492          & 0.9621          & 0.9874          & 0.4800          & 0.5993          & 0.6133          \\
    \midrule[0.3pt]
    WAN~2.1 (25)           & 2.00$\times$ & 59646          & 0.9434          & 0.9597          & \textbf{0.9879} & 0.4250          & 0.5893          & 0.5706          \\
    \textbf{Ours ($S=2$)} & \textbf{2.03}$\times$ & \textbf{57417} & \textbf{0.9462} & \textbf{0.9598} & 0.9859          & \textbf{0.4950} & \textbf{0.5975} & \textbf{0.6114} \\
    \midrule[0.3pt]
    \textbf{Ours ($S=3$)} & 2.54$\times$ & 45953 & 0.9299          & 0.9499          & 0.9815          & 0.5700 & 0.5696          & 0.5990 \\
    \bottomrule
  \end{tabular*}
  \caption{\textbf{Quantitative comparison on latent-space video generation on WAN~2.1~\citep{wan}.} Video resolution is maintained at 720P. Our training-free approach demonstrates more than $2\times$ speedup while maintaining high generation quality and outperforming the 25-step full-resolution baseline.}
  \label{tab:wan_vbench_720p}
\end{table*}




% \renewcommand{\arraystretch}{1.0}
% \setlength{\aboverulesep}{0.2ex}
% \begin{table*}[!t]
%   \centering
%   \scriptsize
%   \begin{tabular*}{\textwidth}{@{\extracolsep{\fill}} l ccccccc @{}}
%     \toprule
%     \makecell{Method} &
%     \makecell{Speedup $\uparrow$} &
%     \makecell{Subject\\Consistency $\uparrow$} &
%     \makecell{Background\\Consistency $\uparrow$} &
%     \makecell{Motion\\Smoothness $\uparrow$} &
%     \makecell{Dynamic\\Degree $\uparrow$} &
%     \makecell{Aesthetic\\Quality $\uparrow$} &
%     \makecell{Image\\Quality $\uparrow$} \\
%     \midrule
%     WAN2.1 (50 steps)           & 1.00$\times$ & 0.9492          & 0.9621          & 0.9874 & 0.4800          & 0.5993          & 0.6133          \\
%     \midrule[0.3pt]
%     WAN2.1 (25 steps)           & 2.00$\times$ & 0.9434          & 0.9597          & \textbf{0.9879} & 0.4250          & 0.5893          & 0.5706          \\
%     \textbf{Ours ($S=2$)} & 2.03$\times$ & \textbf{0.9462} & \textbf{0.9598} & 0.9859 & \textbf{0.4950} & \textbf{0.5975} & \textbf{0.6114} \\
%     \midrule[0.3pt]
%     \textbf{Ours ($S=3$)} & 2.54$\times$ & 0.9299          & 0.9499          & 0.9815 & 0.5700 & 0.5696          & 0.5990 \\
%     \bottomrule
%   \end{tabular*}
%   \vspace{3pt}
%   \caption{\textbf{Quantitative comparison on latent-space video generation on WAN2.1~\citep{wan}.} Video resolution is maintained at 720P. Our training-free approach demonstrates more than $2\times$ speedup while maintaining high generation quality and outperforming the 25-step full-resolution baseline.}
%   \label{tab:wan_vbench_720p}
% \end{table*}


\subsection{Spectral-transformation-based Fine-tuning}
\label{sec:method:lora}
While Spectral Progressive Diffusion is compatible with pretrained image and video generation models in a training-free manner (Sec.~\ref{sec:method:noise-injection}), it implicitly assumes multi-resolution generation capability of pretrained models across scales $s_{1:S}$. Moreover, our optimal resolution schedule (Sec.~\ref{sec:method:schedule:signal}) may not accurately reflect pretrained models' training dynamics, potentially introducing training-inference gaps. 
In this subsection, we outline a spectral-transformation-based fine-tuning approach that directly follows our Spectral Progressive Diffusion framework and our optimal resolution schedule. We include additional implementation details of our fine-tuning framework in Sec.~\ref{sec:exp:image} and Appendix Sec.~\ref{subsec:image_pixel}.

\paragraph{Resolution-specific velocity targets.} 
We assume $S$ progressive resolution scales $s_{1:S}$ and their corresponding $\delta$-optimal resolution transition times $t_{1:S-1}$ from Sec.~\ref{sec:method:schedule:signal} (denoted here as $t_i$ instead of $t^*_i$ for simplicity). For resolution stage $i$, the model operates at scale $s_i$ for $t \in (t_i, \tilde{t}_{i-1}]$, where $\tilde{t}_{i-1}$ is the aligned timestep obtained from Eq.~\eqref{eq:time-align} after expanding from $s_{i-1}$ to $s_i$ at transition time $t_{i-1}$. The state at the beginning of the stage $i$ is hence the enlarged and timestep-aligned $\tilde{\mathbf{x}}_{\tilde{t}_{i - 1}}^{s_{i}}$. The model is fine-tuned to follow a straight-line path from $\tilde{\mathbf{x}}_{\tilde{t}_{i - 1}}^{s_{i}}$ to the standard flow-matching state at the next transition time:
\begin{equation}
  \mathbf{x}_{t_i}^{s_i}
  =
  (1-t_i)\mathbf{x}_0^{s_i}
  +
  t_i\boldsymbol{\epsilon}^{s_i}.
  \label{eq:lora-stagewise}
\end{equation}
where $\mathbf{x}_0^{s_i}$ denotes the clean data sample at scale $s_i$.
The corresponding stage-specific velocity target is therefore:
\begin{equation}
  \mathbf{v}^{s_i}
  =
  \frac{
  \tilde{\mathbf{x}}_{\tilde{t}_{i - 1}}^{s_{i}}
  -
  \mathbf{x}_{t_i}^{s_i}
  }{
  \tilde{t}_{i - 1} - t_i
  }.
  \label{eq:lora-stagewise-velocity}
\end{equation}

\paragraph{Training and inference.}
At each training step we sample $\mathbf{x}_0 \sim p_\text{data}$ and $t \sim \mathcal{U}(0,1)$. We assign $t$ to the resolution stage $i$ satisfying $t \in (t_i,t_{i-1}]$, and then sample $\boldsymbol{\epsilon}^{s_i} \sim \mathcal{N}(0, I)$ at scale $s_i$. We construct the stage input $\tilde{\mathbf{x}}_{\tilde{t}_{i - 1}}^{s_{i}}$ by
applying the same spectral noise expansion and timestep alignment used at
inference, where newly introduced spectral coefficients are filled using
$t_{i-1} \cdot T_\Phi(\boldsymbol{\epsilon}^{s_i})$. We use the same
$\boldsymbol{\epsilon}^{s_i}$ to construct the endpoint
$\mathbf{x}^{s_i}_{t_i}$ from Eq.~\eqref{eq:lora-stagewise}, so that $\tilde{\mathbf{x}}_{\tilde{t}_{i - 1}}^{s_{i}}$ and $\mathbf{x}^{s_i}_{t_i}$ are correlated through a shared noise realization. The training sample is the point on the straight path between them:
\begin{equation}
  \mathbf{x}_t^{s_i}
  =
  \mathbf{x}_{t_i}^{s_i}
  +
  (t-t_i)\,\mathbf{v}^{s_i},
\end{equation}
with target velocity $\mathbf{v}^{s_i}$ from
Eq.~\eqref{eq:lora-stagewise-velocity}.

We fine-tune a pretrained model $\mathbf{v}_\theta$ by minimizing
\begin{equation}
  \mathcal{L}(\theta) \;=\; \mathbb{E}_{t,\, \mathbf{x}_0,\, \boldsymbol{\epsilon}}\!\left[\,\|\,\mathbf{v}_\theta(\mathbf{x}_t^{s_i}, t) - \mathbf{v}^{s_i}\,\|^2\,\right].
  \label{eq:lora-loss}
\end{equation}
At inference, we use the same procedure as in Sec.~\ref{sec:method:noise-injection} with the fine-tuned model, using the same resolution schedule $t_{1:S-1}$ and transform $T_\Phi$.
